% !PRINTED 488
\weberpage{488}
Aus (10) folgt noch weiter:
\weq{12}{f_2(\omega)f_2\!\left(\dfrac{\omega}{n}\right)=\sqrt2\,e^{-\frac{r\pi i}{12}},}
\weq{13}{f(\omega)f\!\left(\dfrac{\omega}{n}\right)f_1(\omega)f_1\!\left(\dfrac{\omega}{n}\right)=\sqrt2\,e^{\frac{r\pi i}{12}}.}
Wenn wir $n=23$ und $r=0$ annehmen, so gibt die Gleichung \S\,75, (8)
mit Benutzung von (12):
\weq{14}{f(\omega)f\!\left(\dfrac{\omega}{n}\right)-f_1(\omega)f_1\!\left(\dfrac{\omega}{n}\right)=2+\sqrt2,}
woraus durch zweimaliges Quadrieren mit Benutzung von (11) und (13) folgt:
\[
x^3+\frac{1}{x^3}=36+26\,\sqrt2;
\]
hieraus leitet man die einfachere Gleichung ab:
\weq{15}{x+\dfrac{1}{x}=3+\sqrt2,\quad \sqrt2\,x=f_1^2(\sqrt{-46}).}
3. Wir machen endlich noch eine Anwendung der Transformationsgleichungen
des \S\,76 für einen zusammengesetzten Transformationsgrad auf die
Determinante $-39$.

Setzen wir
\[
u=f(\sqrt{-39}),\quad v=f\!\left(\sqrt{\dfrac{-13}{3}}\right),\quad \frac{u^2+v^2}{uv}=z,
\]
so ist in der auf $n=39$ bezüglichen Gleichung (22), \S\,76 zu setzen:
\[
A=\frac{u^2}{v^2}-2\,\frac{v}{u},\quad B=\frac{v^2}{u^2}-2\,\frac{u}{v},
\]
und die erwähnte Gleichung gibt:
\[
z^3+z^2-5\,z-6=0,
\]
woraus nach Abwerfung des Faktors $z+2$ die folgende hervorgeht:
\weq{16}{z^2-z-3=0,\quad z=\dfrac{1+\sqrt{13}}{2}.}
Zwischen $u$ und $v$ besteht aber andererseits eine
Transformationsgleichung dritter Ordnung (\S\,73):
\[
\frac{u^6v^6-8}{u^3v^3}=\frac{u^{12}+v^{12}}{u^6v^6}=z^6-6\,z^4+9\,z^2-2,
\]
was sich mit Hilfe von (16) auf die Form $z^4-2=\dfrac{27+7\sqrt{13}}{2}$
bringen läßt.
